\documentclass[11pt]{combiArticle}
\usepackage{tikz,amsmath,amsfonts,latexsym}
\usepackage{amssymb}
\usetikzlibrary{shapes,arrows,calc}
\usepackage{calc}
\usepackage{parskip}
\usepackage{enumitem}
\setlist{itemsep=2pt, topsep=2pt, parsep=0pt}

\definecolor{myblue}{RGB}{0, 84, 159}

\lang{de}

\begin{document}
\course{Schichtplan Manager}
\assistent{Linus Palm}
\contact{palm@combi.rwth-aachen.de}
\title{Modell -- Kurzfassung}
\date{09. November 2026}
\maketitle

Der Schichtplan Manager erstellt für einen Monat und eine Station automatisch einen Dienstplan. In der Demo
ist das die fiktive Musterstation. Dabei hält er feste Regeln immer ein und verfolgt mehrere Ziele in einer festen Rangfolge.

\section{Eingaben}

\begin{itemize}
  \item \textbf{Monat}, für den geplant wird. Die Rechenzeit ist fest: \textbf{30 Sekunden}.
  \item \textbf{Mitarbeitende} mit ihrer Qualifikation: 9 Fachkräfte (drei davon in Teilzeit mit 80 bis 90\,\%),
        eine Hilfskraft und zwei Azubis.
  \item \textbf{Sollzeit und Gutschriften} je Person: 38,5 Stunden je Woche, in Teilzeit anteilig; Urlaub und
        Fortbildung werden mit dem Tagessoll gutgeschrieben.
  \item \textbf{Schichten} mit Uhrzeiten und anrechenbaren Stunden: Früh (F), Zwischen (Z), Spät (S), Nacht (N).
  \item \textbf{Mindestbesetzung} je Tag, Schicht und Qualifikation, jeweils 0 bis 5 Personen.
  \item \textbf{Einschränkungen}: Urlaub, Fortbildung, freier Tag, nicht verfügbar, „nur bestimmte Schichten“.
  \item \textbf{Wünsche}: freier Tag, bestimmte Schicht frei, Wunschtag oder Wunschschicht.
  \item Optional der gespeicherte \textbf{Dienstplan des Vormonats}, damit die Regeln auch über den Monatswechsel
        gelten. Er muss unverändert, genau der Vormonat und angenommen sein. Ohne ihn gelten die Tage vor dem Monat
        als frei.
\end{itemize}

\section{Feste Regeln}

Jeder erstellte Dienstplan hält diese Regeln ein:
\begin{itemize}
  \item Höchstens \textbf{ein Dienst pro Person und Tag}; der Dienst zählt mit der Qualifikation der Person.
  \item \textbf{Einschränkungen} gelten immer. Vor einem Abwesenheitstag gibt es keinen Nachtdienst, am Tag vor
        Urlaubsbeginn keinen Spätdienst.
  \item Höchstens \textbf{10 Arbeitstage in Folge}; Fortbildungstage zählen mit.
  \item \textbf{Nachtdienste} kommen in Blöcken von 2 bis 3 Nächten. Danach folgen 2 freie Tage und am dritten Tag
        kein Frühdienst. Zwischen zwei Nachtblöcken liegen mindestens 7 Tage.
  \item \textbf{Wochenenden}: Samstag und Sonntag dieselbe Schicht oder beide frei; nie zwei Wochenenden in Folge,
        auch nicht über den Monatswechsel. Ein Wochenende gilt als gearbeitet, wenn am Samstag oder Sonntag ein
        Dienst beginnt; ein Nachtdienst am Freitag allein zählt nicht.
  \item Dienste des Vormonats werden \textbf{nicht verändert}, zählen aber bei den Regeln mit. Was dort schon
        feststeht, erzwingt im neuen Monat nichts: Nach einem gearbeiteten Samstag hat der Sonntag am Ersten dieselbe
        Schicht oder ist frei, nach einem freien Samstag ist er offen.
\end{itemize}

Mindestbesetzung, Sollzeit und der Wechsel vom Spät- in den Frühdienst sind bewusst \emph{keine} festen
Regeln, sondern Ziele. So entsteht auch bei Personalmangel ein Plan, und was fehlt, wird als Lücke ausgewiesen.

\section{Ziele in ihrer Rangfolge}

Die Ziele werden nacheinander verbessert. Ein späteres Ziel darf ein früheres nicht verschlechtern; nur bei den
Zielen c), d) und f) ist ein kleiner Spielraum erlaubt, damit die folgenden Ziele noch Wahlmöglichkeiten haben.
\begin{enumerate}
  \item \textbf{Möglichst wenige Lücken}, also unbesetzte Pflichtstellen.
  \item \textbf{Sollzeit einhalten}: Am Monatsende liegt jede Person höchstens 7:40 Stunden über oder unter ihrem Soll.
  \item \textbf{Lücken gleichmäßig verteilen}: Bleiben Lücken, fallen sie gemessen am Bedarf gleichmäßig auf die Tage.
  \item \textbf{Gleichmäßig besetzte Tage}: Jeder Tag ist im gleichen Verhältnis zu seiner Mindestbesetzung besetzt.
  \item \textbf{Wünsche erfüllen}.
  \item \textbf{Gleichmäßig besetzte Schichten}: Zusätzliches Personal eines Tages verteilt sich im Verhältnis
        der Mindestbesetzung auf die Schichten.
  \item \textbf{Möglichst wenige Wechsel vom Spät- in den Frühdienst}.
  \item \textbf{Möglichst oft dieselbe Schicht wie am Vortag}.
\end{enumerate}

Ist die Rechenzeit abgelaufen, wird der beste bis dahin gefundene Plan ausgegeben.

\section{Ausgabe}

\begin{itemize}
  \item \textbf{Der Dienstplan}: für jede Person und jeden Tag die Schicht.
  \item \textbf{Besetzung und Lücken}: je Tag und Schicht die Besetzung im Vergleich zur Mindestbesetzung
        (Ist/Soll) und welche Pflichtstellen mit welcher Qualifikation offen bleiben.
  \item \textbf{Monatskonten}: Soll, Gutschriften, geplante Stunden und Saldo je Person.
  \item \textbf{Unabhängige Prüfung} des Plans gegen gesetzliche und betriebliche Regeln, z.\,B. 11 Stunden Ruhezeit
        zwischen zwei Diensten. Ergebnis: \emph{angenommen}, \emph{abgelehnt} (mit Person, Datum und Regel) oder
        \emph{unvollständig geprüft}. Nur ein angenommener Plan kann als Vormonat für den nächsten Monat
        geladen werden.
  \item \textbf{Kennzahlen}: Lücken, wie viele Monatssalden außerhalb von $\pm$\,7:40 Stunden liegen und die größte
        Abweichung, erfüllte Wünsche und Wechsel vom Spät- in den Frühdienst.
  \item Der Plan lässt sich als \textbf{eine Datei speichern} („Dienstplan \textless Monat\textgreater{}
        \textless Jahr\textgreater{} - Version \textless n\textgreater.json“) und später wieder laden.
\end{itemize}

\end{document}
